% Nouveau EGA V (cote 25): archive cote 25, site page 60 (PDF page 61).
% typescript. Draft compared: none.
% AI-transcribed by Claude Opus 5.5 (claude-opus-5-5, adaptive thinking, effort high), Message Batches API, batch ega25b (msgbatch_01A7n5YuXubUPNK1zsepi8RG), 2026-11-13; not yet checked by eye.
\scanpage{25/60}
\deleted{\illegible{2}} \mspage{55}

Pour comprendre la signification de ce théorème, il conviendrait
\inserted{\deleted{\illegible{5}}}
de noter dès le début du nº \hand{(}que  si S=Spec(k)\hand{)}, alors pour une droite
L ``assez générale'' dans $\hand{\check{\ }}P$ ,  la condition de régularité est vérifiée \inserted{(\unsure{notamment pour S.3.|notamment pour 5.3.})}
(cf ra\xout{t}trappages signalés au  nº\inserted{\unsure{l'avant}}\deleted{suivant}).  Par suite, dans les construc%
tions de ``bons'' pinceaux linéaires  indiquées au début du nº, on pouvait
imposer que les pinceaux envisagés \xout{\unsure{étaient}} satisfaisaient ladite conditi-
on  (ce qui \xout{ne} est une condition\inserted[typed]{de même type mais\hand{,}} distincte de celle qui consiste à
exiger que pour tout \hand{$\zeta\in$}L,  $\phi_{\hand{\zeta}}$  soit $X_{k(\hand{\zeta})}$-régulier ). Il fallait inclu
re la condition en question dans les énoncés récapitulatifs préconisés.
\xout{D'autre part, on peut noter}  D'autre part, pratiquement  10.2. est
utilisé surtout dans la situation 9.5. ,  ce qui  rend  désirable de
n'énoncer la réformulation de  9.5.  en termes de pinceaux qu'après
10.2. , pour pouvoir  inclure dans l'énoncé en question  également l'iso
morphisme du pinceau avec un éclaté (i.e. donner une descr\xout{u}iption de la
situation  permettant une référence commode).  On obtient ainsi un
moyen, pour tout schéma projectif lisse \inserted[typed]{\hand{(}géom \xout{\illegible{1}} conn  X  de di\xout{p}m$\geq$2 \hand{)}}sur un corps infini k,  de trou\xout{x}
ver un\xout{e} sous-schéma fermé lisse\xout{ partout} \inserted[typed]{non vide} de codimension \deleted{\illegible{1}}2 en chacun de
\xout{puisse être}
ses points \xout{(ce qui n'exclut pas qu'il soit vide) et il implique au}
\xout{\unsure{contraire si dim X}} tel que le schéma éclaté admette une fibration
sur P$^1$ , dont toutes les\hand{/}fibres \inserted[typed]{\hand{(}sont géom intègres, et  dont toutes les fibres \hand{)}}sont lisses sont au plus un nombre fini,
ces dernières ayant au plus un point \xout{non lisse (et} géom. singulier, et
un tel point \xout{est} \inserted[typed]{étant} rationnel sur k et géom. singulier ordinaire .  Cela
explique l'importance  de l'étude approfondie (\deleted{à peine} \inserted{tout juste} abordée à l'heure
actuelle) de telles fibrations à fibres singulières, pour réduire dans \xout{x}
une certaine mesure l'étude des variétés projectives lisses de dim d
à celle des(familles à un paramètre de) variétés projectives de dim d-1\hand{,}
\hand{pouvant avoir des singularités ordinaires.}

\margin{60}
